\documentclass[11pt,a4paper]{article}
\usepackage[margin=2.2cm]{geometry}
\usepackage{amsmath,amssymb,booktabs,longtable,microtype}
\usepackage[hidelinks]{hyperref}

\title{Plant Model: Full Nonlinear Formulation\\and its Linear Approximation}
\author{Synthetic Fuel Production --- \texttt{sfp}}
\date{17 September 2026}

\newcommand{\sclip}{\operatorname{sclip}}
\newcommand{\sstep}{\operatorname{sstep}}
\newcommand{\gate}{\operatorname{g}}

\begin{document}
\maketitle

\noindent
Section~\ref{sec:full} is the model integrated by the simulator and differentiated
by the inner NMPC --- one set of equations, used by both, so there is no second
copy to drift. Section~\ref{sec:lp} is the reduced linear model the outer
dispatch layer optimises. Parameter provenance is tagged in
\texttt{docs/ASSUMPTIONS.md}; values quoted here are the defaults in
\texttt{sfp/models/params/*.yaml}. Two thirds of them are \texttt{assumed}.

\section{Full nonlinear model}
\label{sec:full}

\subsection{Conventions}

The plant state is $x\in\mathbb R^{16}$, the control $u\in\mathbb R^{11}$, and
weather $w$ carries plane-of-array irradiance $G$ ($\mathrm{W/m^2}$), air temperature
$T_a$ ($^\circ$C), wind speed $v$ ($\mathrm{m/s}$), relative humidity and
pressure. Each machine takes a \emph{(setpoint, enable)} pair; the battery takes
two powers and the array a curtailment fraction.

\textbf{Sign convention: positive electrical power is power \emph{consumed} from
the DC bus.} The array is therefore a negative load and battery discharge is
negative.

Three smoothing primitives appear throughout, all from
\texttt{sfp/models/mathx.py}. They exist because the inner layer differentiates
these expressions and a hard switch has no usable gradient:
\begin{align}
\sstep(z;\varepsilon) &= \tfrac12\bigl(1+\tanh(z/\varepsilon)\bigr), &
\sclip(z;a,b,\varepsilon) &\approx \min(\max(z,a),b).
\end{align}
Every machine is gated by $\gate(u) = \sstep(u_{\text{enable}}-\tfrac12;\,0.05)$,
which is $2.06\times10^{-9}$ when off and $1$ when on. Integration is fixed-step
RK4 with a zero-order hold on $u$ and $w$.

\subsection{PV array}

No states. Input: curtailment fraction $\gamma\in[0,1]$.

Faiman (2008) cell temperature, linear temperature derating about STC, constant
inverter efficiency and a hard clip at the AC rating:
\begin{align}
T_{\text{cell}} &= T_a + \frac{G}{u_0 + u_1 v}, &
f_{\text{derate}} &= (1-s)(1-\ell)(1-d)^{y},\\
P_{\text{dc}} &= P_{\text{dc}}^{\text{rated}}\,\frac{G}{G_{\text{ref}}}
  \bigl(1+\gamma_P (T_{\text{cell}}-T_{\text{ref}})\bigr) f_{\text{derate}}, &
P_{\text{ac}} &= \min\!\bigl(\eta_{\text{inv}} P_{\text{dc}},\, P_{\text{ac}}^{\text{rated}}\bigr),
\end{align}
with $P_{\text{ac}}^{\text{rated}} = P_{\text{dc}}^{\text{rated}}/\text{DC:AC}$
and delivered power $(1-\gamma)P_{\text{ac}}$. Clipping loss and deliberate
curtailment are reported separately --- the brief asks for curtailment, and off
grid it is simply PV that cannot be absorbed.

\begin{center}\small
\begin{tabular}{llr}
\toprule
$\gamma_P$ & power temperature coefficient & $-0.0035\ \mathrm{K^{-1}}$\\
$T_{\text{ref}}, G_{\text{ref}}$ & STC reference & $25\,^\circ$C, $1000\ \mathrm{W/m^2}$\\
$u_0, u_1$ & Faiman coefficients & $25\ \mathrm{W/(m^2K)}$, $6.84\ \mathrm{Ws/(m^3K)}$\\
$\eta_{\text{inv}}$, DC:AC & inverter & $0.975$, $1.25$\\
$s,\ \ell,\ d$ & soiling, wiring, annual degradation & $0.03,\ 0.02,\ 0.005\ \mathrm{yr^{-1}}$\\
\bottomrule
\end{tabular}
\end{center}

\subsection{Battery}

States: state of charge $\sigma$, fractional capacity fade $\phi$, cumulative
equivalent full cycles $n$. Inputs: charge and discharge power, both
non-negative.
\begin{align}
E_{\text{use}} &= E_{\text{nom}}\max(1-\phi,\,10^{-3}), &
\dot\sigma &= \frac{\eta_c P_{\text{ch}} - P_{\text{dis}}/\eta_d}{E_{\text{use}}}
   - \frac{\kappa_{\text{sd}}}{86400}\,\sigma,\\
\dot n &= \frac{P_{\text{ch}}+P_{\text{dis}}}{2E_{\text{nom}}}, &
\dot\phi &= \frac{\phi_{\text{eol}}}{N_{\text{cyc}}}\,S(\sigma)\,\dot n
   + \frac{\phi_{\text{cal}}}{T_{\text{yr}}} .
\end{align}

Charge and discharge are \emph{separate non-negative} variables rather than one
signed power. Nothing forbids both being positive at once; what prevents it is
that the wear term makes it strictly worse than either alone, which is how this
formulation avoids an explicit complementarity constraint.

The depth-of-discharge stress factor $S$ encodes
$N(\mathrm{DoD}) = N_{\text{ref}}(\mathrm{DoD}_{\text{ref}}/\mathrm{DoD})^{k}$:
\begin{equation}
S(\sigma) = C\,\Bigl[\max\bigl(\underbrace{2|\sigma-\bar\sigma|/\mathrm{DoD}_{\text{ref}}}_{\text{depth}},\ f_{\min}\bigr)\Bigr]^{k-1},
\qquad \bar\sigma=\tfrac12(\sigma_{\max}+\sigma_{\min}),
\end{equation}
with $C$ chosen so a cycle at the rated depth costs exactly one rated cycle.
\textbf{This is a local proxy for a path-dependent quantity}: depth of discharge
is properly a property of a \emph{cycle}, recovered by rainflow counting, and a
differential model cannot see a cycle. Without it, throughput is all that
matters and a large pack swinging gently is charged the same per kWh as a small
one cycling to its limits --- which flatters small packs precisely in a sizing
sweep.

\begin{center}\small
\begin{tabular}{llr}
\toprule
$\eta_c,\eta_d$ & one-way efficiencies & $0.96$, $0.96$\\
$\sigma_{\min},\sigma_{\max}$ & SoC window & $0.10$, $0.95$\\
$\kappa_{\text{sd}}$ & self-discharge & $0.002\ \mathrm{d^{-1}}$\\
$N_{\text{cyc}}$, $\mathrm{DoD}_{\text{ref}}$ & rated cycle life & $6000$ at $0.80$\\
$k$, $f_{\min}$ & DoD exponent, stress floor & $3.0$, $0.25$\\
$\phi_{\text{eol}}$, $\phi_{\text{cal}}$ & end-of-life fade, calendar fade & $0.20$, $0.015\ \mathrm{yr^{-1}}$\\
\bottomrule
\end{tabular}
\end{center}

\subsection{Air contactor}

No states --- the sorbent it loads lives in \S\ref{sub:solids}. Input: airflow
fraction $\varphi$ and enable.

Capture and fan power scale oppositely with airflow, which is the whole point of
the subsystem:
\begin{align}
V &= \varphi\,\gate(u)\,V_{\max}, &
\rho &= \frac{p}{R T_a}, &
\mathrm{NTU} &= \mathrm{NTU}_{\text{d}}\frac{V_{\text{d}}}{\max(V,10^{-3})},\\
\eta_{\text{cap}} &= 1-e^{-\mathrm{NTU}}, &
f_T &= \exp\!\Bigl[-\frac{E_a}{R}\Bigl(\frac1{T_a}-\frac1{T_{\text{ref}}}\Bigr)\Bigr], &
f_{\mathrm{RH}} &= \frac{\mathrm{rh}/(\mathrm{rh}+k_h)}{\mathrm{rh}_{\text{ref}}/(\mathrm{rh}_{\text{ref}}+k_h)},
\end{align}
\begin{equation}
r_{\text{carb}} = V\rho\,y_{\mathrm{CO_2}}\,\eta_{\text{cap}}\,f_T\,f_{\mathrm{RH}}\,f_{\text{load}},
\qquad
P_{\text{fan}} = P_{\text{fan}}^{\text{rated}}\varphi^{3} + P_{\text{idle}}\gate(u).
\end{equation}
$\mathrm{NTU}_{\text{d}} = -\ln(1-\eta_{\text{d}})$; residence time falls as
$1/V$ in a fixed bed, so $\mathrm{NTU}\propto 1/V$. The load taper
$f_{\text{load}} = \sclip\bigl((1-L)/(1-L_{\text{taper}})\bigr)$ stops the fan
pushing CO$_2$ into a saturated bed.

\textbf{Diminishing returns are severe.} At reference sizing, $20^\circ$C and
60\,\% RH: $60\,\%$ flow captures $0.233$\,mol/s for $4.7$\,kW
($128\ \mathrm{kWh/tCO_2}$); $100\,\%$ flow captures $0.279$\,mol/s for
$20.4$\,kW ($462\ \mathrm{kWh/tCO_2}$). Twenty percent more CO$_2$ for four and
a third times the fan power --- a power-follow controller that dumps surplus
solar into the fans burns most of it for nothing.

\begin{center}\small
\begin{tabular}{llr}
\toprule
$V_{\max}$, $V_{\text{d}}$ & max and design airflow & $40$, $24\ \mathrm{m^3/s}$\\
$\eta_{\text{d}}$ & design single-pass capture & $0.65$\\
$P_{\text{fan}}^{\text{rated}}$, $P_{\text{idle}}$ & fan power & $20$\,kW, $0.4$\,kW\\
$y_{\mathrm{CO_2}}$ & ambient mole fraction & $420$\,ppm (measured)\\
$E_a$, $k_h$, $L_{\text{taper}}$ & activation, RH half-sat., taper & $20$\,kJ/mol, $0.3$, $0.9$\\
\bottomrule
\end{tabular}
\end{center}

\subsection{Sorbent inventory}
\label{sub:solids}

States: $n_{\mathrm{CaO}}$, $n_{\mathrm{CaCO_3}}$, mean cycle number $\bar N$.
No inputs --- it is driven by the contactor and calciner rates.
\begin{align}
\dot n_{\mathrm{CaO}} &= r_{\text{calc}} - r_{\text{carb}} + \dot m - \pi r_{\text{calc}}, &
\dot n_{\mathrm{CaCO_3}} &= r_{\text{carb}} - r_{\text{calc}}, &
\dot{\bar N} &= \frac{r_{\text{calc}} - \dot m \bar N}{n_{\text{tot}}} .
\end{align}
One full pass of the inventory through the kiln advances $\bar N$ by exactly one;
make-up dilutes the mean age. Grasa--Abanades deactivation sets the capacity:
\begin{equation}
X(\bar N) = \frac{1}{k\bar N + (1-X_r)^{-1}} + X_r,
\qquad X_r = 0.075,\ k = 0.52 .
\end{equation}
$\bar N$ is kept \emph{continuous} rather than integer specifically so it stays
differentiable inside the NLP. Because $\mathrm{d}X/\mathrm{d}\bar N<0$, every
calcination permanently shrinks the capacity of the \emph{whole} inventory ---
the marginal cost the controller pays per mole calcined, and the reason it is
reluctant to cycle the loop harder than the methane is worth.

$n_{\text{tot}} = 50\,000$\,mol; make-up and purge are both zero by default, so
deactivation is irreversible in the runs reported.

\subsection{Calciner}

State: kiln temperature $T$. Inputs: heater fraction and enable.

Three independent limits multiply --- thermodynamics, kinetics, feedstock:
\begin{align}
p_{\text{eq}}(T) &= A_B e^{-B_B/T} \quad\text{(Baker 1962)}, &
f_{\text{drive}} &= \sclip\!\Bigl(1-\frac{p_{\text{op}}}{p_{\text{eq}}}\Bigr), \\
f_{\text{kin}} &= A_k e^{-E_a/(RT)}, &
f_{\text{feed}} &= \sclip\!\Bigl(\frac{n_{\mathrm{CaCO_3}}}{n_{\text{ref}}}\Bigr),
\end{align}
\begin{equation}
r_{\text{calc}} = r_{\max}\,f_{\text{drive}}\,\sclip(f_{\text{kin}})\,f_{\text{feed}}\,\gate(u).
\end{equation}
Energy balance, with all four terms in watts:
\begin{equation}
C\,\dot T = \underbrace{f_h P_h^{\text{rated}}\eta_h}_{\text{heater}}
 - \underbrace{(UA + c_v v)(T-T_a)}_{\text{loss}}
 - \underbrace{r_{\text{calc}}\,\Delta H_{\text{calc}}}_{\text{reaction}}
 - \underbrace{r_{\text{calc}} M_{\mathrm{CaCO_3}} c_p \max(T-T_{\text{feed}},0)}_{\text{sensible}} .
\end{equation}
The heater fraction carries an over-temperature interlock. \textbf{The kiln's
thermal inertia is the plant's commitment problem}: $C = 3\times10^6$\,J/K
against a 150\,kW heater means \textbf{5.6 hours from ambient to 1200\,K}, so a
kiln worth lighting is a kiln worth running for a while. That single fact is why
the outer layer keeps a kiln temperature state at all.

\begin{center}\small
\begin{tabular}{llr}
\toprule
$C$, $P_h^{\text{rated}}$, $\eta_h$ & thermal capacity, heater & $3.0\times10^6$\,J/K, $150$\,kW, $0.95$\\
$UA$, $c_v$ & loss coefficient, wind term & $17\ \mathrm{W/K}$, $0.6\ \mathrm{W/(K\,m/s)}$\\
$A_B$, $B_B$ & Baker equilibrium & $4.137\times10^7$\,atm, $20474$\,K\\
$p_{\text{op}}$ & kiln operating pressure & $0.3$\,atm\\
$A_k$, $E_a$ & Arrhenius & $2.207\times10^5\ \mathrm{s^{-1}}$, $120$\,kJ/mol\\
$r_{\max}$, $T_{\max}$ & feeder throughput, limit & $0.45$\,mol/s, $1273$\,K\\
$\Delta H_{\text{calc}}$ & endothermic & $+178$\,kJ/mol\\
\bottomrule
\end{tabular}
\end{center}

\subsection{PEM electrolyser}

States: stack temperature $T$, voltage degradation $\delta V$. Input: current
fraction and enable. $i = \varphi\,\gate(u)\,i_{\max}$.
\begin{align}
V_{\text{rev}}(T) &= V_{\text{rev}}^0 + \alpha_T (T-298.15), &
\eta_{\text{act}} &= \frac{RT}{\alpha_c F}\,\operatorname{asinh}\!\Bigl(\frac{i}{2i_0}\Bigr),\\
\mathrm{ASR}(T) &= \mathrm{ASR}_0\bigl(1+\beta(T-T_{\text{set}})\bigr), &
V_{\text{cell}} &= V_{\text{rev}} + \eta_{\text{act}} + i\,\mathrm{ASR} + \delta V .
\end{align}
The activation term is Butler--Volmer in $\operatorname{asinh}$ form, which is
exact and---unlike the Tafel approximation---smooth down to zero current.
\begin{align}
r_{\mathrm{H_2}} &= \frac{\eta_F N_c i A}{2F}, &
P_{\text{stack}} &= N_c V_{\text{cell}} i A, &
Q_{\text{gen}} &= N_c i A (V_{\text{cell}} - V_{\text{tn}}),\\
P_{\text{total}} &= P_{\text{stack}} + P_{\text{aux}}\gate(u) + \kappa Q_{\text{gen}}, &
C\dot T &= Q_{\text{gen}} - Q_{\text{cool}} - UA(T-T_a), &
\dot{\delta V} &= \delta_r \varphi .
\end{align}
Heat is written as the voltage difference $V_{\text{cell}}-V_{\text{tn}}$ rather
than (electrical in $-$ chemical out), to avoid differencing two large, nearly
equal numbers.

\textbf{Efficiency peaks at part load and then falls}, because $\eta_{\text{act}}$
grows like $\log i$ while the ohmic term grows like $i$. ``Run flat out when
sunny'' is therefore provably suboptimal --- one of the few places where the
economics and the physics point the same way.

\begin{center}\small
\begin{tabular}{llr}
\toprule
$N_c$, $A$, $i_{\max}$ & cells, area, max current density & $105$, $1000\ \mathrm{cm^2}$, $1.5\ \mathrm{A/cm^2}$\\
$i_{\min}$ & minimum load (gas crossover) & $0.10\,i_{\max}$\\
$V_{\text{rev}}^0$, $V_{\text{tn}}$, $\alpha_T$ & reversible, thermoneutral & $1.229$\,V, $1.481$\,V, $-0.9$\,mV/K\\
$i_0$, $\alpha_c$, $\mathrm{ASR}_0$, $\beta$ & kinetics, resistance & $10^{-4}\ \mathrm{A/cm^2}$, $0.5$, $0.15\ \Omega\mathrm{cm^2}$, $-0.005\ \mathrm{K^{-1}}$\\
$\eta_F$, $P_{\text{aux}}$, $\kappa$ & Faraday eff., auxiliaries, parasitic & $0.98$, $15$\,kW, $0.02$\\
$C$, $UA$, $T_{\max}$ & thermal & $2.5\times10^5$\,J/K, $150$\,W/K, $353.15$\,K\\
\bottomrule
\end{tabular}
\end{center}

\subsection{Gas buffers and water balance}

Hydrogen and CO$_2$ tanks, no inputs:
\begin{equation}
\dot n_{\mathrm{H_2}} = a_{\mathrm{H_2}} r_{\text{ely}} - 4 r_{\text{sab}},
\qquad
\dot n_{\mathrm{CO_2}} = a_{\mathrm{CO_2}} r_{\text{calc}} - r_{\text{sab}},
\end{equation}
where $a_\bullet\in[0,1]$ are smooth acceptance gates that close as a tank fills.
Production the tanks cannot accept is \emph{vented and recorded}, not silently
dropped from the mass balance: a power-follow controller can vent nearly half the
hydrogen it makes, having paid ${\sim}56$\,kWh/kg for it, purely because it never
looked at the tank level.

Water, with condensate recycle $\phi_r = 0.95$:
\begin{equation}
\dot m_{w} = \dot m_{\text{makeup}} + 2\phi_r M_{\mathrm{H_2O}} r_{\text{sab}} - M_{\mathrm{H_2O}} r_{\text{ely}} .
\end{equation}
Net consumption is ${\approx}2.25$\,kg H$_2$O per kg CH$_4$ after recycle --- a
real constraint at an arid site. Capacities: $30\,000$\,mol H$_2$,
$6000$\,mol CO$_2$, $5000$\,kg water; CO$_2$ compression costs $17.5$\,kJ/mol.

\subsection{Sabatier reactor}

States: temperature $T$, catalyst activity $a$, cumulative methane $m$. Input:
feed fraction $\varphi$ and enable.
\begin{align}
X_{\text{eq}}(T) &= \frac{1}{1+\exp\bigl((T-T_{50})/w\bigr)}, &
\mathrm{Da} &= \mathrm{Da}_{\text{ref}}\,e^{-\frac{E_a}{R}(\frac1T-\frac1{T_{\text{ref}}})}\,\frac{a}{\max(\varphi,0.02)},\\
X &= X_{\text{eq}}\bigl(1-e^{-\mathrm{Da}}\bigr), &
r_{\text{sab}} &= \varphi F_{\max} X \, \Lambda(w).
\end{align}
$X_{\text{eq}}$ \emph{falls} with temperature because the reaction is
exothermic; that is the term that stops the controller running the reactor as
hot as the catalyst allows. Residence time rises as the feed is turned down, so
$\mathrm{Da}$ rises --- conversion improves at part load.

The availability factor is the hard stoichiometric coupling:
\begin{equation}
\Lambda = \min\!\Bigl(\sclip\bigl(\tfrac{n_{\mathrm{CO_2}}}{50}\bigr),\,
                      \sclip\bigl(\tfrac{n_{\mathrm{H_2}}}{200}\bigr)\Bigr)\cdot\sstep(m_w-1).
\end{equation}
\textbf{Whichever of CO$_2$ or H$_2$ runs short throttles the reactor, no matter
how much of the other is banked.} This is what forces the two upstream chains to
be co-scheduled rather than optimised separately, and it is the single most
important coupling in the plant.
\begin{align}
C\dot T &= r_{\text{sab}}(-\Delta H) + Q_{\text{pre}} - Q_{\text{cool}} - UA(T-T_a) - Q_{\text{sens}}, \\
\dot a &= -k_d \exp\!\Bigl(\frac{T-T_{\text{ref}}^d}{T_{\text{scale}}}\Bigr) a,
\qquad \dot m = r_{\text{sab}} M_{\mathrm{CH_4}} .
\end{align}
Deactivation is exponential in temperature, so thermal excursions are
permanently expensive.

\begin{center}\small
\begin{tabular}{llr}
\toprule
$F_{\max}$ & max CO$_2$ feed & $0.12$\,mol/s\\
$T_{50}$, $w$ & equilibrium logistic & $810$\,K, $70$\,K\\
$\mathrm{Da}_{\text{ref}}$, $E_a$ & kinetics at $573$\,K & $6.0$, $80$\,kJ/mol\\
$T_{\text{ign}}$, $T_{\text{target}}$, $T_{\max}$ & light-off, target, sintering & $523$, $573$, $823$\,K\\
$k_d$, $T_{\text{scale}}$ & deactivation & $10^{-8}\ \mathrm{s^{-1}}$, $50$\,K\\
$\Delta H$ & exothermic & $-165$\,kJ/mol\\
\bottomrule
\end{tabular}
\end{center}

\subsection{DC bus}

The subsystems are coupled only through the power balance, evaluated in
registration order so that each publishes what the next needs:
\begin{equation}
\underbrace{P_{\text{fan}} + P_{\text{kiln}} + P_{\text{ely}} + P_{\text{sab}}
 + w_{\text{comp}} r_{\text{calc}}}_{\text{loads}}
 + \underbrace{P_{\text{ch}} - P_{\text{dis}}}_{\text{battery}}
 - \underbrace{(1-\gamma)P_{\text{ac}}}_{\text{array}} = 0 .
\end{equation}
A regulatory layer below the controller enforces this in the real plant: it caps
PV and battery to what the hardware can deliver and, when demand still exceeds
supply, \textbf{sheds load in a fixed priority order} --- calciner, contactor,
electrolyser, Sabatier --- recording every intervention and every undervoltage
trip. A controller that needs the bus to save it is failing.

\section{Linear approximation for the dispatch LP}
\label{sec:lp}

The outer layer is a MILP over $n = 240$ hourly intervals: 7447 columns, 6964
rows, 24 integers. It carries \emph{all} the economics and, apart from the kiln,
no dynamics. See \texttt{docs/DISPATCH\_NMPC.md} \S2 for the row census.

\subsection{What is dropped, and what that costs}

\begin{center}\small
\begin{tabular}{lll}
\toprule
\textbf{Full model} & \textbf{LP treatment} & \textbf{Consequence}\\
\midrule
16 states & 7 states & fast thermal states invisible\\
$\eta_{\text{cap}}(V)$, $\eta_{\text{act}}(i)$, $X(T)$ & concave PWL in dispatch & exact relaxation (\S\ref{sub:pwl})\\
reactor $T$, electrolyser $T$ & dropped & light-off and warm-up unmodelled\\
kiln $T$ & \textbf{kept} --- balance is linear & 5.6\,h warm-up preserved\\
$r_{\text{calc}}(T)$ Arrhenius$\times$Baker & linear bound + binary & genuinely non-convex (\S\ref{sub:kiln})\\
catalyst, stack degradation & dropped & no thermal-cycling cost\\
hourly grid & --- & \textbf{sub-hourly buffering invisible}\\
\bottomrule
\end{tabular}
\end{center}

The last line is not a rounding error. Measured closed-loop, the LP's planned
operating profit is flat across a $15\times$ range of battery capacity while
realised profit rises $2.5\times$: an hourly average of supply and demand simply
does not contain the deficits a battery exists to cover, so the layer cannot
price storage. See \texttt{experiments/BATTERY\_SIZING.md}.

\subsection{States and controls}

$z_k\in\mathbb R^7 = (\sigma,\ n_{\mathrm{CaCO_3}},\ \bar N,\ n_{\mathrm{H_2}},\
n_{\mathrm{CO_2}},\ m_w,\ T_{\text{kiln}})$; 24 columns per interval: four
commitments $e_j$, seven PWL segment fills, kiln power/rate/indicator, four start
indicators, two battery powers, curtailment, two vent rates and make-up water.

\subsection{Piecewise-linear rate maps}
\label{sub:pwl}

For the contactor, electrolyser and reactor, each nonlinear rate curve is
replaced by segment fills $\delta_{j,m}\in[0,w_{j,m}]$:
\begin{equation}
d_j = \sum_m \delta_{j,m}, \qquad r_j = \sum_m a_{j,m}\delta_{j,m},
\qquad \delta_{j,m} \le w_{j,m} e_j ,
\end{equation}
with slopes $a_{j,m}$ \textbf{strictly decreasing}. That concavity is what makes
the relaxation \emph{exact}: an optimiser maximising output fills the steepest
segment first unprompted, so no ordering binaries are needed. The builder raises
if a sampled curve is not concave --- without it the formulation is silently
wrong rather than merely approximate. Electrical draw is
$P_j = \text{idle}_j e_j + d_j$ for the power-dispatched machines.

\subsection{Kiln: the one irreducible binary}
\label{sub:kiln}

The energy balance is linear exactly as written in the full model, with the
sensible-heat term evaluated at the operating temperature rather than at $T_k$
--- the only approximation, worth a few percent of duty during warm-up:
\begin{equation}
C\frac{T_{k+1}-T_k}{\Delta t} = \eta P^{\text{kiln}}_k - UA(T_k - T_a) - \Delta H\,r^{\text{cal}}_k .
\end{equation}
The \emph{rate} is the problem. The set
$0 \le r \le \max\bigl(0,\ \alpha(T-T_{\text{on}})\bigr)$ is genuinely
non-convex, and writing the cap directly rearranges to
$T \ge T_{\text{on}} + r/\alpha$, which at $r=0$ demands a permanently hot kiln.
An indicator is unavoidable:
\begin{equation}
r \le \bar r y, \qquad
r \le \alpha(T-T_{\text{on}}) + M(1-y), \qquad
y \le e, \qquad M = \alpha(T_{\text{on}}-250).
\end{equation}
Relaxing $y$ to $[0,1]$ permits
$r^\star = \bar r(\alpha(T-T_{\text{on}})+M)/(\bar r+M)$, which at 841\,K allows
$0.26$\,mol/s where the real kiln calcines nothing --- and reproduces exactly the
failure it was added to prevent: calcination scheduled on a cold kiln, sorbent
saturation, production collapsing to 17\,kg/day. $y$ is therefore binary over the
first 24 hours and relaxed beyond, where only the terminal inventory depends on
it.

\subsection{Dynamics}

Seven equalities per interval, all linear:
\begin{align}
\sigma_{k+1} &= \sigma_k\Bigl(1-\tfrac{\Delta t\,\kappa_{\text{sd}}}{86400}\Bigr)
  + \tfrac{\Delta t\,\eta_c P^{\text{ch}}_k}{E} - \tfrac{\Delta t\,P^{\text{dis}}_k}{\eta_d E}, &
n^{\mathrm{CaCO_3}}_{k+1} &= n^{\mathrm{CaCO_3}}_k + \Delta t(r^{\text{cont}}_k - r^{\text{cal}}_k),\\
\bar N_{k+1} &= \bar N_k + \Delta t\,r^{\text{cal}}_k/n_{\text{tot}}, &
n^{\mathrm{H_2}}_{k+1} &= n^{\mathrm{H_2}}_k + \Delta t(r^{\text{ely}}_k - 4r^{\text{sab}}_k - v^{\mathrm{H_2}}_k),\\
n^{\mathrm{CO_2}}_{k+1} &= n^{\mathrm{CO_2}}_k + \Delta t(r^{\text{cal}}_k - r^{\text{sab}}_k - v^{\mathrm{CO_2}}_k), &
m^w_{k+1} &= m^w_k + \Delta t(\dot m^w_k + 2\phi_r M r^{\text{sab}}_k - M r^{\text{ely}}_k).
\end{align}

\subsection{Remaining constraints}

\textbf{Bus balance} (the block whose dual is $\lambda$), written
demand-minus-supply so relaxing it means energy from nowhere and the dual is a
positive price:
$\sum_j P_{j,k} + w_{\text{comp}} r^{\text{cal}}_k + P^{\text{ch}}_k - P^{\text{dis}}_k = (1-\gamma_k)\mathrm{PV}_k$.

\textbf{Reactor availability}, a linear envelope of the plant's $\min(\cdot,\cdot)$
taper --- needing no indicator, because the taper reaches zero exactly \emph{at}
the box floor so $r^{\text{sab}}=0$ stays feasible:
\begin{equation}
r^{\text{sab}}_k \le \bar r_{\text{sab}}\frac{n^{\mathrm{CO_2}}_k-\underline n_{\mathrm{CO_2}}}{50},
\qquad
r^{\text{sab}}_k \le \bar r_{\text{sab}}\frac{n^{\mathrm{H_2}}_k-\underline n_{\mathrm{H_2}}}{200}.
\end{equation}

\textbf{Commitment}: minimum electrolyser load
$\sum_m \delta_{\text{ely},m,k}\ge \underline d\,e_k$; start counting
$s_{j,k}\ge e_{j,k}-e_{j,k-1}$, with interval 0 measured against what is
\emph{already running} rather than against zero; and minimum up-time
$\tau_j s_{j,k}\le\sum_{l=k}^{k+\tau_j-1}e_{j,l}$ with $\tau=3$ for the calciner
and $2$ for the reactor.

\subsection{Objective}

\begin{equation}
\min \sum_{k}\Bigl[-p_{\mathrm{CH_4}}M_{\mathrm{CH_4}}\Delta t\,r^{\text{sab}}_k
+ \frac{c_{\text{sorb}}\Delta t\,r^{\text{cal}}_k}{n_{\text{tot}}}
+ \frac{c_{\text{batt}}\Delta t}{2E}(P^{\text{ch}}_k+P^{\text{dis}}_k)
+ c_w\Delta t\,\dot m^w_k
+ \sum_j c^{\text{start}}_j s_{j,k}\Bigr] - V(z_n),
\end{equation}
with start costs $0.50/40/5/40$\,EUR for contactor/calciner/electrolyser/reactor
and terminal value $V$ pricing leftover inventory at the methane it would become.
Without $V$ the layer empties every buffer at the horizon end, which is locally
optimal and globally wrong.

$c_{\text{batt}} = \mathrm{CAPEX}/(E\cdot\mathrm{DoD}_{\text{ref}}\cdot N_{\text{cyc}}\cdot\eta_{\text{rt}})$
$= 0.0565$\,EUR/kWh at the reference pack, against a night-time shadow price of
${\approx}0.043$\,EUR/kWh. \textbf{$p_{\mathrm{CH_4}}$ is a control tuning
parameter, not a market price}: set below the crossover the optimiser declines
to cycle the battery, curtails a third of the array, runs the pack flat and trips
the bus. It is fixed at $5.50$\,EUR/kg, above the crossover with margin.

\end{document}
